\section{Proposed Method: Trust-Enhanced Graph-Based Recommendation Model (GTERM)}
The rapid expansion of digital platforms has resulted in an abundant accumulation of user-generated data, underscoring the demand for efficient recommendation systems capable of filtering this information to deliver personalized content. Graph Neural Networks (GNNs) have gained prominence as an effective tool for these systems, particularly due to their ability to model complex user-item interactions through graph structures \cite{wu2019neural, kipf2017semi}. GNNs excel in capturing both connectivity and complex interaction patterns, thus providing a robust framework for enhancing recommendations through structural insights. Despite these advancements, GNNs often neglect the critical role of trust dynamics in recommendation scenarios \cite{velickovic2018graph}.

Research on trust mechanisms in recommendation systems has highlighted their potential to enhance recommendation reliability by incorporating user trust scores \cite{massa2007trust, guo2017trust}. These trust-based approaches aim to mitigate biases from unreliable interactions and increase transparency by elucidating the rationale behind recommendations. Nevertheless, many trust-based methodologies operate independently of sophisticated network-based techniques, missing the opportunity to integrate the comprehensive interaction modeling capabilities of GNNs \cite{konstas2009social, ma2011recommender}. The incorporation of trust into GNN frameworks remains inadequately explored, presenting an opportunity for innovation that our proposed model seeks to exploit.

The integration of trust and GNNs faces several challenges. Conventional GNN-based recommendation models frequently lack explicit trust metrics, limiting their capacity to deliver interpretable and user-focused recommendations. Specifically, while GNNs adeptly model intricate interaction structures, they may not distinguish interactions based on their trustworthiness. Moreover, the opacity of such models often compromises user confidence and acceptance \cite{zhang2014explicit}. Addressing these issues demands a methodology that effectively embeds trust into the recommendation process while enhancing interpretability and trust in the system.

In response to these challenges, we introduce the Trust-Enhanced Graph-Based Recommendation Model (GTERM), which directly incorporates trust metrics into the GNN framework for recommendation applications. GTERM converts raw interaction data into a trust-augmented graph, forming the foundation for modeling user-item interactions in a trust-sensitive context. Utilizing graph convolutional and attention mechanisms, GTERM accurately captures intricate interaction patterns, emphasizing trust-enriched interactions to improve recommendation accuracy and transparency without substantially altering the GNN architecture. This approach addresses existing challenges while ensuring recommendations remain both data-driven and trust-aware, thereby enhancing user acceptance and understanding.

\begin{itemize}
    \item A novel data preprocessing technique is developed for constructing trust-infused graphs from raw user-item interactions, facilitating more robust and intricate interaction modeling.
    \item A GNN-based interaction modeling approach is designed to incorporate trust metrics, refining feature propagation and interaction predictions with an emphasis on trustworthy connections.
    \item A trust-enhanced recommendation framework is introduced, which integrates graph-based predictions with trust data, providing transparent and explainable recommendations that improve user engagement.
    \item GTERM is empirically validated on diverse datasets, demonstrating notable improvements in accuracy and transparency over existing methods, with comprehensive ablation studies conducted to emphasize the role of trust integration in enhancing system performance.
\end{itemize}
